\section{引言}
\label{sec:intro}

张量网络（tensor network）模型长期以来一直是建模复杂量子系统的最先进方法~\citep{white1992, fannes1992, orus2019}，但直到最近才被用作机器学习模型~\citep{novikov2015, cohen2016, stoudenmire2016, novikov2017, han2018, stoudenmire2018, cheng2019}。与神经网络不同，张量网络不使用非线性激活函数，而是依靠乘性相互作用来捕捉数据中的复杂关联。这赋予张量网络一种便于证明一般性理论结果的数学结构，例如几乎所有深层张量网络与其浅层对应物在表达能力上的分离~\citep{cohen2016}。然而，这些独特的数学性质尚未被用于发展新的\emph{操作性}（operational）能力，而这类能力将为张量网络模型在现实机器学习任务中获得更广泛的采用提供更多实际理由。

在这项工作中，我们将一种循环张量网络——\emph{均匀矩阵乘积态}（uniform matrix product state, u-MPS）——应用于概率序列建模（probabilistic sequence modeling）任务，并指出了 u-MPS 在评估与生成能力方面的若干新颖特性。尽管 u-MPS 具有循环性质，我们证明其序列输入可以高度并行地处理，长度为 $n$ 的序列可以在 $\bigO(\log n)$ 的并行时间内完成评估。虽然深层循环神经网络（RNN）难以并行化，这一困难此前曾推动人们为序列处理任务开发非循环架构（如 \citep{gehring2017, vaswani2017}），但我们的发现表明，循环张量网络是获得更高并行性的另一条途径。

我们进一步证明，u-MPS 模型具有与正则表达式（regex）结构紧密相关的惊人生成能力。标准的自回归（autoregressive）模型只能以流式方式在给定提示条件下生成序列，而我们发现 u-MPS 能够从由条件正则表达式 $R$ 定义的各种各样的分布中采样。我们的采样算法可以高效地从 u-MPS 学到的概率分布中产生无偏样本，条件是输出序列匹配给定的正则表达式 $R$。标准自回归采样对应于取 $R = p\Sigma^*$（其中 $p$ 是前缀字符串，$\Sigma^*$ 是匹配所有序列的正则表达式），其他特殊情形还包括填空采样（$R = p\Sigma^*s$，其中 $s$ 为后缀），以及生成被约束为包含某个目标短语 $t$ 的样本（$R=\Sigma^*t\Sigma^*$）。

除了允许生成具有丰富内部结构的序列之外，这些技术还实现了新颖的 regularization 形式：在训练过程中可以对 u-MPS 模型进行惩罚或激励，使其生成匹配某个目标模式的字符串。这种 regularization 可以应用于多种具有挑战性的任务，包括自动代码生成以及缓解语言模型中的性别偏见。在合成数据集与真实结构化文本数据集上的实验证实了这些新颖的并行性、采样与 regularization 优势，并表明 u-MPS 能够成功地将短字符串中存在的非局域关联泛化到长度显著更大的序列上。

\paragraph{贡献总结}
我们首次实现了用于概率序列建模的 u-MPS，并揭示了该模型的若干卓越能力\footnote{产生本文结果的代码可在 \url{https://github.com/jemisjoky/umps_code} 获取。}。u-MPS 中非线性激活函数的缺失使我们能够在训练与推理阶段利用一种并行评估方法。我们发展了将 u-MPS「转移算符」（transfer operators）的结构与正则表达式结构联系起来的新技术，这进而催生了灵活的递归采样算法以及针对 u-MPS 的新颖 regularization 形式。我们期望这些技术能为序列生成模型的设计开辟重要的新研究方向，其中语言建模是一个尤其有前景的领域。

\paragraph{相关工作}
张量网络在机器学习中的若干重要先前应用包括压缩大型神经网络权重~\citep{novikov2015}、证明深层网络与浅层网络表达能力的分离~\citep{cohen2016}，以及用于有监督~\citep{stoudenmire2016, novikov2017, glasser2018} 与无监督（unsupervised）~\citep{han2018, stoudenmire2018, cheng2019} 学习任务。与之尤为相关的是~\citep{stokes2019}，其中（非均匀的）MPS 借助密度矩阵重正化群（DMRG）算法被训练为定长二进制序列的生成模型。人们还提出了多种多样的张量网络架构，作为建模与理解自然语言的理论工具，例如~\citep{pestun2017, coecke2010, gallego2017, degiuli2019}。我们采样算法中所使用的完全正映射与隐量子马尔可夫模型（HQMM）~\citep{monras2010,srinivasan2018} 中使用的映射类似，同样可以借助量子信息理论的概念加以诠释。

这项工作可视为~\citep{pestun2017a} 的延续，该文从理论视角将 u-MPS 作为语言模型引入，但没有给出本文所提供的并行化、采样或实验结果。我们的采样算法是~\citep{han2018} 中引入的定长 Born 机器（Born machine）算法的重大推广（后者又沿承了~\citep{ferris2012} 的方法），并且凭借 u-MPS 的循环性质，能够生成长度任意的离散序列。u-MPS 模型等价于二次加权有限自动机~\citep{bailly2011}，并在较小程度上与范数可观测算子模型（NOOM）~\citep{zhao2010} 相关，它也是线性（二阶）RNN 的一个实例~\citep{rabusseau2019}。线性（一阶）RNN 在并行化与可解释性方面的优势曾在~\citep{martin2018, foerster2017input} 中被描述。

与先前工作的一个关键区别在于，我们为评估 regex 结构的概率分布以及从其采样发展了一般性技术，据我们所知这些技术是全新的。这些技术不仅适用于 u-MPS，也适用于具有类似内部结构的一大类模型，例如加权有限自动机（WFA）~\citep{droste2009}、隐马尔可夫模型（HMM）~\citep{rabiner1986} 与预测状态表示~\citep{littman2002}。因此，我们期望本文针对 regex 采样、regularization 与并行评估所发展的算法能够立即推广到这些参数化有效概率分布的模型中的任何一个。
